% ==============================================================================
% BLOQUE 3: ANGELO PRIETO (Slides 17 a 24 — 8 diapositivas)
% ==============================================================================

% SLIDE 17: PLANTA DE 3ER ORDEN
\section{Compensación de Sistemas}
\subsection{Planta de 3er Orden}
\begin{frame}
  \frametitle{Planta de Tercer Orden y Parámetros}

  \begin{columns}[t, onlytextwidth]
    \begin{column}{0.48\textwidth}
      \begin{tarjeta}
        {\fontsize{9pt}{11pt}\selectfont \textbf{\color{NavyDark} Función de Transferencia:}}\\[2pt]
        {\fontsize{8pt}{10pt}\selectfont
        \[ G(s) = \frac{K}{\left(1 + \frac{s}{p_1}\right)\left(1 + \frac{s}{p_2}\right)\left(1 + \frac{s}{p_3}\right)} \]
        \begin{itemize}\setlength{\itemsep}{2.5pt}
          \item \textbf{Ganancia estática en continua:}
          $K = 900 \implies |G(0)| = \mathbf{\qty{59.08}{\decibel}}$.
          \item \textbf{Ancho de banda base ($BW_{\text{LA}}$):}
          $BW = f_{-3\text{dB}} \approx \mathbf{\qty{11.59}{\kilo\hertz}}$.
          \item Dominado por el primer polo de $12\text{ kHz}$.
        \end{itemize}
        }
      \end{tarjeta}
    \end{column}

    \begin{column}{0.48\textwidth}
      \begin{tarjeta}
        {\fontsize{9pt}{11pt}\selectfont \textbf{\color{NavyDark} Polos Asignados (Grupo 5):}}\\[2pt]
        {\fontsize{8pt}{10pt}\selectfont
        \begin{itemize}\setlength{\itemsep}{3pt}
          \item \textbf{Primer polo ($f_{p1}$):}
          $f_{p1} = \qty{12}{\kilo\hertz} \implies p_1 \approx \qty{75398}{\radian\per\second}$
          \item \textbf{Segundo polo ($f_{p2}$):}
          $f_{p2} = \qty{70}{\kilo\hertz} \implies p_2 \approx \qty{439823}{\radian\per\second}$
          \item \textbf{Tercer polo ($f_{p3}$):}
          $f_{p3} = \qty{140}{\kilo\hertz} \implies p_3 \approx \qty{879646}{\radian\per\second}$
          \item Planta de orden 3 con alta ganancia y polos en alta frecuencia.
        \end{itemize}
        }
      \end{tarjeta}
    \end{column}
  \end{columns}
\end{frame}

% SLIDE 18: ANÁLISIS DE BODE
\subsection{Respuesta en Frecuencia}
\begin{frame}
  \frametitle{Respuesta en Frecuencia a Lazo Abierto (Bode)}

  \begin{columns}[c, onlytextwidth]
    \begin{column}{0.58\textwidth}
      \centering
      \adjustbox{max width=\linewidth, max height=5.5cm, keepaspectratio}{\begin{tikzpicture}
    \pgfplotsset{
        bode plot/.style={
            /pgf/number format/read comma as period,
            grid=both,
            major grid style={solid, gray!45},
            minor grid style={dotted, gray!35},
            axis line style={black, line width=0.6pt},
            axis x line* = bottom,
            axis y line* = left,
            extra description/.code={
                \draw[gray!45] (rel axis cs:0,1) -- (rel axis cs:1,1) -- (rel axis cs:1,0);
            },
            width=15.5cm,
            height=6.5cm,
            xmin=100, xmax=10000000,
            enlarge x limits=false,
            tick label style={font=\small},
            label style={font=\normalsize},
        }
    }

    % Subplot 1: Magnitud
    \begin{semilogxaxis}[
        bode plot,
        ylabel={Magnitud [\si{\decibel}]},
        ymin=-70, ymax=75,
        ytick={-60, -40, -20, 0, 20, 40, 60},
        xticklabel=\empty,
        legend pos=south west,
        legend style={font=\small, fill=white, fill opacity=0.9, text opacity=1},
        name=magplot,
    ]
        % Curva exacta
        \addplot[line width=1.3pt, blue!85!black] table [col sep=comma, x=frequency, y=mag_exact] {sources/plots/bode_planta.csv};
        \addlegendentry{$G(s)$}

        % Nivel 0 dB
        \addplot[black!70, densely dotted, line width=0.8pt] coordinates {(100, 0) (10000000, 0)};
        \node[black!70, anchor=south east, font=\scriptsize] at (axis cs:9000000, 1.5) {$0\,\text{dB}$};

        % Ancho de banda a -3 dB
        \draw[teal!80!black, dashed, line width=1.0pt] (axis cs:11593, -70) -- (axis cs:11593, 56.07);
        \addplot[only marks, mark=*, mark size=2.5pt, teal!80!black] coordinates {(11593, 56.07)};
        \node[teal!80!black, fill=white, fill opacity=0.9, text opacity=1, inner sep=2pt, font=\small, anchor=south] 
            at (axis cs:11593, 62) {$BW = 11,59\,\text{kHz}$};
        \draw[teal!80!black, -stealth, line width=0.8pt] (axis cs:11593, 61.5) -- (axis cs:11593, 56.7);

        % Guía vertical f_45 (MF = 45°)
        \draw[purple, dashed, line width=1.0pt] (axis cs:52901, -70) -- (axis cs:52901, 43.44);
        \addplot[only marks, mark=*, mark size=2.6pt, purple] coordinates {(52901, 43.44)};
        \node[purple, fill=white, fill opacity=0.9, text opacity=1, inner sep=1.5pt, font=\small, anchor=south] 
            at (axis cs:52901, 45.5) {$43,44\,\text{dB}$};

        % Guía vertical f_180
        \draw[magenta!80!black, dashed, line width=1.0pt] (axis cs:110993, -70) -- (axis cs:110993, 32.14);
        \addplot[only marks, mark=square*, mark size=2.4pt, magenta!80!black] coordinates {(110993, 32.14)};
        \node[magenta!80!black, fill=white, fill opacity=0.9, text opacity=1, inner sep=1.5pt, font=\small, anchor=north east] 
            at (axis cs:105000, 30.5) {$32,14\,\text{dB}$};

        % Cota de Margen de Ganancia para MF = 45°
        \draw[orange!90!black, dotted, line width=0.8pt] (axis cs:52901, 43.44) -- (axis cs:230000, 43.44);
        \draw[orange!90!black, dotted, line width=0.8pt] (axis cs:110993, 32.14) -- (axis cs:230000, 32.14);
        \draw[stealth-stealth, orange!90!black, line width=1.2pt] 
            (axis cs:200000, 32.14) -- (axis cs:200000, 43.44)
            node[midway, right=2pt, font=\small, fill=white, fill opacity=0.9, text opacity=1, inner sep=1.5pt] {$MG_{45} = 11,30\,\text{dB}$};
    \end{semilogxaxis}

    % Subplot 2: Fase (posicionado estrictamente DEBAJO de magplot)
    \begin{semilogxaxis}[
        bode plot,
        xlabel={$f$ [\si{\hertz}]},
        ylabel={Fase [$^\circ$]},
        ymin=-280, ymax=10,
        ytick={-270, -225, -180, -135, -90, -45, 0},
        at={(magplot.south west)},
        anchor=north west,
        yshift=-0.7cm,
    ]
        % Curva exacta
        \addplot[line width=1.3pt, blue!85!black] table [col sep=comma, x=frequency, y=phase_exact] {sources/plots/bode_planta.csv};

        % Línea de referencia -135° (MF = 45°)
        \addplot[purple, dashed, line width=1.0pt] coordinates {(100, -135) (10000000, -135)};
        \node[purple, fill=white, fill opacity=0.85, text opacity=1, anchor=south east, font=\small, inner sep=1.5pt] 
            at (axis cs:9000000, -132) {$-135^\circ\;(MF=45^\circ)$};

        % Línea de referencia -180°
        \addplot[red!80!black, dotted, line width=1.0pt] coordinates {(100, -180) (10000000, -180)};
        \node[red!80!black, fill=white, fill opacity=0.85, text opacity=1, anchor=south east, font=\small, inner sep=1.5pt] 
            at (axis cs:9000000, -177) {$-180^\circ$};

        % Cruce f_45
        \draw[purple, dashed, line width=1.0pt] (axis cs:52901, -280) -- (axis cs:52901, -135);
        \addplot[only marks, mark=*, mark size=2.6pt, purple] coordinates {(52901, -135)};
        \node[purple, fill=white, fill opacity=0.9, text opacity=1, inner sep=2pt, font=\small, anchor=north] 
            at (axis cs:52901, -143) {$f_{45} = 52,90\,\text{kHz}$};

        % Cruce f_180
        \draw[magenta!80!black, dashed, line width=1.0pt] (axis cs:110993, -280) -- (axis cs:110993, -180);
        \addplot[only marks, mark=square*, mark size=2.4pt, magenta!80!black] coordinates {(110993, -180)};
        \node[magenta!80!black, fill=white, fill opacity=0.9, text opacity=1, inner sep=2pt, font=\small, anchor=north] 
            at (axis cs:110993, -188) {$f_{180} = 110,99\,\text{kHz}$};
    \end{semilogxaxis}
\end{tikzpicture}
}
    \end{column}

    \begin{column}{0.40\textwidth}
      \begin{tarjeta}
        {\fontsize{8.5pt}{10pt}\selectfont \textbf{\color{NavyDark} Comportamiento Asintótico:}}\\[2pt]
        {\fontsize{7pt}{8.5pt}\selectfont
        \begin{itemize}\setlength{\itemsep}{2pt}
          \item \textbf{$f < 12\text{ kHz}$:} Ganancia $+59.08\text{ dB}$, fase $\ang{0}$.
          \item \textbf{$12\text{k} < f < 70\text{ kHz}$:} Actúa $p_1$; pendiente $-20\ \text{dB/década}$, fase hacia $\ang{-90}$.
          \item \textbf{$70\text{k} < f < 140\text{ kHz}$:} Actúan $p_1, p_2$; pendiente $-40\ \text{dB/década}$, cruce por $\ang{-180}$.
          \item \textbf{$f > 140\text{ kHz}$:} Actúan 3 polos; pendiente $-60\ \text{dB/década}$, fase hacia $\ang{-270}$.
        \end{itemize}
        }
      \end{tarjeta}
    \end{column}
  \end{columns}
\end{frame}

% SLIDE 19: DIAGNÓSTICO DE INESTABILIDAD
\subsection{Diagnóstico de Inestabilidad}
\begin{frame}
  \frametitle{Diagnóstico de Inestabilidad a Lazo Cerrado}

  \begin{columns}[t, onlytextwidth]
    \begin{column}{0.48\textwidth}
      \begin{tarjeta}[colframe=AlertRed]
        {\fontsize{9pt}{11pt}\selectfont \textbf{\color{AlertRed} Diagnóstico Numérico Directo:}}\\[2pt]
        {\fontsize{8pt}{10pt}\selectfont
        \begin{itemize}\setlength{\itemsep}{2.5pt}
          \item \textbf{Cruce de fase por $\ang{-180}$:}
          $f_{180} \approx \mathbf{\qty{110.99}{\kilo\hertz}}$
          \item \textbf{Ganancia en el cruce de fase:}
          $|G(f_{180})| \approx \mathbf{\qty{32.14}{\decibel}}$ ($40.5$ veces)
          \item \textbf{Margen de Ganancia directo:}
          $MG = 0 - 32.14 = \mathbf{-32.14\ \unit{\decibel}}$
          \item \textbf{Cruce de ganancia por $\qty{0}{\decibel}$:}
          $f_{0\si{\decibel}} \approx \qty{464.42}{\kilo\hertz}$ con fase $\ang{-243.17}$ ($MF = \ang{-63.17}$).
        \end{itemize}
        }
      \end{tarjeta}
    \end{column}

    \begin{column}{0.48\textwidth}
      \begin{tarjeta}
        {\fontsize{9pt}{11pt}\selectfont \textbf{\color{NavyDark} Consecuencia Física:}}\\[2pt]
        {\fontsize{8pt}{10pt}\selectfont
        \begin{itemize}\setlength{\itemsep}{2.5pt}
          \item Con realimentación unitaria directa, la planta es \textbf{incondicionalmente inestable}.
          \item En $110.99\text{ kHz}$ el desfasaje de $\ang{-180}$ transforma la realimentación negativa en positiva con amplificación de $40.5$ veces.
          \item Los polos de lazo cerrado migran al semiplano derecho $\to$ auto-oscilación violenta.
          \item \textbf{Solución:} Es mandatorio diseñar redes de compensación en frecuencia.
        \end{itemize}
        }
      \end{tarjeta}
    \end{column}
  \end{columns}
\end{frame}

% SLIDE 20: CALIBRACIÓN MF = 45°
\subsection{Calibración para MF = 45°}
\begin{frame}
  \frametitle{Criterio de Calibración para Margen de Fase $MF = \ang{45}$}

  \begin{columns}[c, onlytextwidth]
    \begin{column}{0.48\textwidth}
      \begin{tarjeta}[colframe=AlertGreen]
        {\fontsize{9pt}{11pt}\selectfont \textbf{\color{AlertGreen} Determinación Gráfica en Bode:}}\\[2pt]
        {\fontsize{8pt}{10pt}\selectfont
        \begin{itemize}\setlength{\itemsep}{2.5pt}
          \item \textbf{Fase requerida para $MF = \ang{45}$:}
          $\phi = \ang{-180} + \ang{45} = \mathbf{\ang{-135}}$
          \item \textbf{Frecuencia para $\phi = \ang{-135}$ ($f_{45}$):}
          $f_{45} \approx \mathbf{\qty{52.90}{\kilo\hertz}}$
          \item \textbf{Ganancia disponible en $f_{45}$:}
          $|G(f_{45})| \approx \mathbf{\qty{43.44}{\decibel}}$ ($148.6$ veces)
          \item \textbf{Margen de Ganancia disponible ($MG_{45}$):}
          $MG_{45} = 43.44 - 32.14 = \mathbf{11.30\ \unit{\decibel}} > \qty{6}{\decibel}$
        \end{itemize}
        }
      \end{tarjeta}
    \end{column}

    \begin{column}{0.48\textwidth}
      \begin{tarjeta}
        {\fontsize{9pt}{11pt}\selectfont \textbf{\color{NavyDark} Significado en Diseño:}}\\[2pt]
        {\fontsize{8pt}{10pt}\selectfont
        \begin{itemize}\setlength{\itemsep}{2.5pt}
          \item Atenuar puramente con un divisor resistivo exigiría $\beta = 1/148.6 = \qty{-43.44}{\decibel}$.
          \item $MG_{45} = \qty{11.30}{\decibel} > \qty{6}{\decibel}$ garantiza que la planta admite operar con $MF = \ang{45}$ manteniendo margen seguro.
          \item Se abordarán dos estrategias pasivas:
          \begin{itemize}\setlength{\itemsep}{1.5pt}
            \item Compensador de Retardo (Lag).
            \item Red de Atraso-Adelanto (Lead-Lag).
          \end{itemize}
        \end{itemize}
        }
      \end{tarjeta}
    \end{column}
  \end{columns}
\end{frame}

% SLIDE 21: RETARDO SÍNTESIS
\subsection{Compensador de Retardo}
\begin{frame}
  \frametitle{Síntesis del Compensador de Retardo (Lag)}

  \begin{columns}[c, onlytextwidth]
    \begin{column}{0.48\textwidth}
      \begin{tarjeta}
        {\fontsize{8.5pt}{10.5pt}\selectfont \textbf{\color{NavyDark} Topología y Dimensionamiento:}}\\[2pt]
        {\fontsize{7.5pt}{9.2pt}\selectfont
        \[ G_c(s) = \frac{1}{1 + sRC} = \frac{1}{1 + \frac{s}{2\pi f_d}}, \quad f_d = \qty{1.2}{\kilo\hertz} \]
        \begin{itemize}\setlength{\itemsep}{2pt}
          \item Capacitor comercial adoptado: $C = \mathbf{\qty{100}{\nano\farad}}$ (film).
          \item Resistencia calculada:
          $R = \frac{1}{2\pi f_d C} \approx \qty{1326.29}{\ohm}$.
          \item \textbf{Normalización comercial ($\qty{1}{\percent}$):}
          $R = \mathbf{\qty{1.33}{\kilo\ohm}}$ (Serie E96).
          \item \textbf{Principio del Polo Dominante:}
          El polo en $1.2\text{ kHz}$ derriba la ganancia antes de que los polos en $70\text{k}$ y $140\text{ kHz}$ aporten desfasaje.
        \end{itemize}
        }
      \end{tarjeta}
    \end{column}

    \begin{column}{0.50\textwidth}
      \centering
      \adjustbox{max width=\linewidth, max height=4.6cm, keepaspectratio}{\begin{circuitikz}[american, scale=1.0]
    \draw (0,2) node[left] {$V_i$} 
          to[short, o-] (0.5,2)
          to[R, l=$R$, -*] (3,2)
          to[short, -o] (4.2,2) node[right] {$V_o$};
    \draw (3,2) 
          to[C, l=$C$] (3,0) node[ground] {};
    \draw (0,0) node[left] {} 
          to[short, o-] (0.5,0) node[ground] {};
    \draw (4.2,0) node[right] {} 
          to[short, o-] (3.7,0) node[ground] {};
\end{circuitikz}
}
    \end{column}
  \end{columns}
\end{frame}

% SLIDE 22: RETARDO BODE
\begin{frame}
  \frametitle{Compensador de Retardo: Respuesta y Pérdida de Ancho de Banda}

  \begin{columns}[c, onlytextwidth]
    \begin{column}{0.58\textwidth}
      \centering
      \adjustbox{max width=\linewidth, max height=5.5cm, keepaspectratio}{\begin{tikzpicture}
    \pgfplotsset{
        bode plot/.style={
            /pgf/number format/read comma as period,
            grid=both,
            major grid style={solid, gray!45},
            minor grid style={dotted, gray!35},
            axis line style={black, line width=0.6pt},
            axis x line* = bottom,
            axis y line* = left,
            extra description/.code={
                \draw[gray!45] (rel axis cs:0,1) -- (rel axis cs:1,1) -- (rel axis cs:1,0);
            },
            width=15.5cm,
            height=6.5cm,
            xmin=100, xmax=10000000,
            enlarge x limits=false,
            tick label style={font=\small},
            label style={font=\normalsize},
        }
    }

    % Subplot 1: Magnitud
    \begin{semilogxaxis}[
        bode plot,
        ylabel={Magnitud [\si{\decibel}]},
        ymin=-100, ymax=75,
        ytick={-80, -60, -40, -20, 0, 20, 40, 60},
        xticklabel=\empty,
        legend pos=south west,
        legend style={font=\small, fill=white, fill opacity=0.9, text opacity=1},
        name=magplot,
    ]
        % Curva exacta
        \addplot[line width=1.3pt, blue!85!black] table [col sep=comma, x=frequency, y=mag_exact] {sources/plots/bode_retardo.csv};
        \addlegendentry{$G_{\text{ret}}(s)$}

        % Nivel 0 dB
        \addplot[black!70, densely dotted, line width=0.8pt] coordinates {(100, 0) (10000000, 0)};
        \node[black!70, anchor=south east, font=\scriptsize] at (axis cs:9000000, 1.5) {$0\,\text{dB}$};

        % Ancho de banda a -3 dB
        \draw[teal!80!black, dashed, line width=1.0pt] (axis cs:1188, -100) -- (axis cs:1188, 56.07);
        \addplot[only marks, mark=*, mark size=2.5pt, teal!80!black] coordinates {(1188, 56.07)};
        \node[teal!80!black, fill=white, fill opacity=0.9, text opacity=1, inner sep=2pt, font=\small, anchor=south] 
            at (axis cs:1188, 62) {$BW = 1,19\,\text{kHz}$};
        \draw[teal!80!black, -stealth, line width=0.8pt] (axis cs:1188, 61.5) -- (axis cs:1188, 56.7);

        % Guía vertical f_45
        \draw[purple, dashed, line width=1.0pt] (axis cs:9962, -100) -- (axis cs:9962, 38.25);
        \addplot[only marks, mark=*, mark size=2.6pt, purple] coordinates {(9962, 38.25)};
        \node[purple, fill=white, fill opacity=0.9, text opacity=1, inner sep=1.5pt, font=\small, anchor=south] 
            at (axis cs:9962, 40.5) {$38,25\,\text{dB}$};

        % Guía vertical f_180
        \draw[magenta!80!black, dashed, line width=1.0pt] (axis cs:24354, -100) -- (axis cs:24354, 25.21);
        \addplot[only marks, mark=square*, mark size=2.4pt, magenta!80!black] coordinates {(24354, 25.21)};
        \node[magenta!80!black, fill=white, fill opacity=0.9, text opacity=1, inner sep=1.5pt, font=\small, anchor=north east] 
            at (axis cs:23000, 23.5) {$25,21\,\text{dB}$};

        % Cota de Margen de Ganancia para MF = 45°
        \draw[orange!90!black, dotted, line width=0.8pt] (axis cs:9962, 38.25) -- (axis cs:55000, 38.25);
        \draw[orange!90!black, dotted, line width=0.8pt] (axis cs:24354, 25.21) -- (axis cs:55000, 25.21);
        \draw[stealth-stealth, orange!90!black, line width=1.2pt] 
            (axis cs:50000, 25.21) -- (axis cs:50000, 38.25)
            node[midway, right=2pt, font=\small, fill=white, fill opacity=0.9, text opacity=1, inner sep=1.5pt] {$MG_{45} = 13,04\,\text{dB}$};
    \end{semilogxaxis}

    % Subplot 2: Fase (posicionado estrictamente DEBAJO de magplot)
    \begin{semilogxaxis}[
        bode plot,
        xlabel={$f$ [\si{\hertz}]},
        ylabel={Fase [$^\circ$]},
        ymin=-370, ymax=10,
        ytick={-360, -315, -270, -225, -180, -135, -90, -45, 0},
        at={(magplot.south west)},
        anchor=north west,
        yshift=-0.7cm,
    ]
        % Curva exacta
        \addplot[line width=1.3pt, blue!85!black] table [col sep=comma, x=frequency, y=phase_exact] {sources/plots/bode_retardo.csv};

        % Línea de referencia -135° (MF = 45°)
        \addplot[purple, dashed, line width=1.0pt] coordinates {(100, -135) (10000000, -135)};
        \node[purple, fill=white, fill opacity=0.85, text opacity=1, anchor=south east, font=\small, inner sep=1.5pt] 
            at (axis cs:9000000, -132) {$-135^\circ\;(MF=45^\circ)$};

        % Línea de referencia -180°
        \addplot[red!80!black, dotted, line width=1.0pt] coordinates {(100, -180) (10000000, -180)};
        \node[red!80!black, fill=white, fill opacity=0.85, text opacity=1, anchor=south east, font=\small, inner sep=1.5pt] 
            at (axis cs:9000000, -177) {$-180^\circ$};

        % Cruce f_45
        \draw[purple, dashed, line width=1.0pt] (axis cs:9962, -370) -- (axis cs:9962, -135);
        \addplot[only marks, mark=*, mark size=2.6pt, purple] coordinates {(9962, -135)};
        \node[purple, fill=white, fill opacity=0.9, text opacity=1, inner sep=2pt, font=\small, anchor=north] 
            at (axis cs:9962, -143) {$f_{45} = 9,96\,\text{kHz}$};

        % Cruce f_180
        \draw[magenta!80!black, dashed, line width=1.0pt] (axis cs:24354, -370) -- (axis cs:24354, -180);
        \addplot[only marks, mark=square*, mark size=2.4pt, magenta!80!black] coordinates {(24354, -180)};
        \node[magenta!80!black, fill=white, fill opacity=0.9, text opacity=1, inner sep=2pt, font=\small, anchor=north] 
            at (axis cs:24354, -188) {$f_{180} = 24,35\,\text{kHz}$};
    \end{semilogxaxis}
\end{tikzpicture}
}
    \end{column}

    \begin{column}{0.40\textwidth}
      \begin{tarjeta}
        {\fontsize{8.5pt}{10pt}\selectfont \textbf{\color{NavyDark} Parámetros (MF = 45°):}}\\[2pt]
        {\fontsize{7pt}{8.5pt}\selectfont
        \begin{itemize}\setlength{\itemsep}{2pt}
          \item \textbf{Colapso de ancho de banda:}\\
          $BW = f_{-3\text{dB}} \approx \mathbf{\qty{1.19}{\kilo\hertz}}$\\
          (pérdida del $\mathbf{89.7\%}$ de banda útil).
          \item \textbf{Frecuencia para $MF = \ang{45}$:}\\
          $f_{45} \approx \qty{9.96}{\kilo\hertz}$ ($|G| \approx \qty{38.25}{\decibel}$).
          \item \textbf{Cruce por $\ang{-180}$:}\\
          $f_{180} \approx \qty{24.35}{\kilo\hertz}$ ($|G| \approx \qty{25.21}{\decibel}$).
          \item \textbf{Margen de Ganancia disponible:}\\
          $MG_{45} = \mathbf{13.04\ \unit{\decibel}} > \qty{6}{\decibel}$.
        \end{itemize}
        }
      \end{tarjeta}
    \end{column}
  \end{columns}
\end{frame}

% SLIDE 23: ATRASO-ADELANTO ESTRUCTURA
\subsection{Compensador de Atraso-Adelanto}
\begin{frame}
  \frametitle{Red de Atraso-Adelanto: Estructura Topológica}

  \begin{columns}[c, onlytextwidth]
    \begin{column}{0.48\textwidth}
      \begin{tarjeta}
        {\fontsize{8.5pt}{10.5pt}\selectfont \textbf{\color{NavyDark} Topología Circuital y Objetivos:}}\\[2pt]
        {\fontsize{7.5pt}{9.2pt}\selectfont
        \begin{itemize}\setlength{\itemsep}{2pt}
          \item \textbf{Estructura:} Rama serie con $R_1$, y rama derivación con $R_2$ en serie con $C$.
          \item \textbf{Función de transferencia de la red:}
          \[ G_{c2}(s) = \frac{1 + s R_2 C}{1 + s (R_1 + R_2) C} = \frac{1 + \frac{s}{2\pi f_z}}{1 + \frac{s}{2\pi f_p}} \]
          \item \textbf{Frecuencias asignadas:}
          Polo $f_p = \qty{220}{\kilo\hertz}$, Cero $f_z = \qty{1.1}{\mega\hertz}$.
          \item \textbf{Preservación de Banda Base:}
          Como $f_p, f_z \gg 12\text{ kHz}$, la banda base útil permanece inalterada.
        \end{itemize}
        }
      \end{tarjeta}
    \end{column}

    \begin{column}{0.50\textwidth}
      \centering
      \adjustbox{max width=\linewidth, max height=4.8cm, keepaspectratio}{\begin{circuitikz}[american, scale=1.0]
    \draw (0,3) node[left] {$V_i$} 
          to[short, o-] (0.5,3)
          to[R, l=$R_1$, -*] (3,3)
          to[short, -o] (4.5,3) node[right] {$V_o$};
    \draw (3,3) 
          to[R, l=$R_2$, *-] (3,1.5)
          to[C, l=$C$] (3,0) node[ground] {};
    \draw (0,0) node[left] {} 
          to[short, o-] (0.5,0) node[ground] {};
    \draw (4.5,0) node[right] {} 
          to[short, o-] (4.0,0) node[ground] {};
\end{circuitikz}
}
    \end{column}
  \end{columns}
\end{frame}

% SLIDE 24: ATRASO-ADELANTO BODE Y RESULTADOS
\begin{frame}
  \frametitle{Red de Atraso-Adelanto: Respuesta y Preservación de Ancho de Banda}

  \begin{columns}[c, onlytextwidth]
    \begin{column}{0.58\textwidth}
      \centering
      \adjustbox{max width=\linewidth, max height=5.5cm, keepaspectratio}{\begin{tikzpicture}
    \pgfplotsset{
        bode plot/.style={
            /pgf/number format/read comma as period,
            grid=both,
            major grid style={solid, gray!45},
            minor grid style={dotted, gray!35},
            axis line style={black, line width=0.6pt},
            axis x line* = bottom,
            axis y line* = left,
            extra description/.code={
                \draw[gray!45] (rel axis cs:0,1) -- (rel axis cs:1,1) -- (rel axis cs:1,0);
            },
            width=15.5cm,
            height=6.5cm,
            xmin=100, xmax=10000000,
            enlarge x limits=false,
            tick label style={font=\small},
            label style={font=\normalsize},
        }
    }

    % Subplot 1: Magnitud
    \begin{semilogxaxis}[
        bode plot,
        ylabel={Magnitud [\si{\decibel}]},
        ymin=-70, ymax=75,
        ytick={-60, -40, -20, 0, 20, 40, 60},
        xticklabel=\empty,
        legend pos=south west,
        legend style={font=\small, fill=white, fill opacity=0.9, text opacity=1},
        name=magplot,
    ]
        % Curva exacta
        \addplot[line width=1.3pt, blue!85!black] table [col sep=comma, x=frequency, y=mag_exact] {sources/plots/bode_atraso_adelanto.csv};
        \addlegendentry{$G_{\text{aa}}(s)$}

        % Nivel 0 dB
        \addplot[black!70, densely dotted, line width=0.8pt] coordinates {(100, 0) (10000000, 0)};
        \node[black!70, anchor=south east, font=\scriptsize] at (axis cs:9000000, 1.5) {$0\,\text{dB}$};

        % Ancho de banda a -3 dB
        \draw[teal!80!black, dashed, line width=1.0pt] (axis cs:11563, -70) -- (axis cs:11563, 56.07);
        \addplot[only marks, mark=*, mark size=2.5pt, teal!80!black] coordinates {(11563, 56.07)};
        \node[teal!80!black, fill=white, fill opacity=0.9, text opacity=1, inner sep=2pt, font=\small, anchor=south] 
            at (axis cs:11563, 62) {$BW = 11,56\,\text{kHz}$};
        \draw[teal!80!black, -stealth, line width=0.8pt] (axis cs:11563, 61.5) -- (axis cs:11563, 56.7);

        % Guía vertical f_45
        \draw[purple, dashed, line width=1.0pt] (axis cs:45109, -70) -- (axis cs:45109, 45.18);
        \addplot[only marks, mark=*, mark size=2.6pt, purple] coordinates {(45109, 45.18)};
        \node[purple, fill=white, fill opacity=0.9, text opacity=1, inner sep=1.5pt, font=\small, anchor=south] 
            at (axis cs:45109, 47) {$45,18\,\text{dB}$};

        % Guía vertical f_180
        \draw[magenta!80!black, dashed, line width=1.0pt] (axis cs:84486, -70) -- (axis cs:84486, 36.22);
        \addplot[only marks, mark=square*, mark size=2.4pt, magenta!80!black] coordinates {(84486, 36.22)};
        \node[magenta!80!black, fill=white, fill opacity=0.9, text opacity=1, inner sep=1.5pt, font=\small, anchor=north east] 
            at (axis cs:80000, 34.5) {$36,22\,\text{dB}$};

        % Cota de Margen de Ganancia para MF = 45°
        \draw[orange!90!black, dotted, line width=0.8pt] (axis cs:45109, 45.18) -- (axis cs:200000, 45.18);
        \draw[orange!90!black, dotted, line width=0.8pt] (axis cs:84486, 36.22) -- (axis cs:200000, 36.22);
        \draw[stealth-stealth, orange!90!black, line width=1.2pt] 
            (axis cs:170000, 36.22) -- (axis cs:170000, 45.18)
            node[midway, right=2pt, font=\small, fill=white, fill opacity=0.9, text opacity=1, inner sep=1.5pt] {$MG_{45} = 8,96\,\text{dB}$};
    \end{semilogxaxis}

    % Subplot 2: Fase (posicionado estrictamente DEBAJO de magplot)
    \begin{semilogxaxis}[
        bode plot,
        xlabel={$f$ [\si{\hertz}]},
        ylabel={Fase [$^\circ$]},
        ymin=-280, ymax=10,
        ytick={-270, -225, -180, -135, -90, -45, 0},
        at={(magplot.south west)},
        anchor=north west,
        yshift=-0.7cm,
    ]
        % Curva exacta
        \addplot[line width=1.3pt, blue!85!black] table [col sep=comma, x=frequency, y=phase_exact] {sources/plots/bode_atraso_adelanto.csv};

        % Línea de referencia -135° (MF = 45°)
        \addplot[purple, dashed, line width=1.0pt] coordinates {(100, -135) (10000000, -135)};
        \node[purple, fill=white, fill opacity=0.85, text opacity=1, anchor=south east, font=\small, inner sep=1.5pt] 
            at (axis cs:9000000, -132) {$-135^\circ\;(MF=45^\circ)$};

        % Línea de referencia -180°
        \addplot[red!80!black, dotted, line width=1.0pt] coordinates {(100, -180) (10000000, -180)};
        \node[red!80!black, fill=white, fill opacity=0.85, text opacity=1, anchor=south east, font=\small, inner sep=1.5pt] 
            at (axis cs:9000000, -177) {$-180^\circ$};

        % Cruce f_45
        \draw[purple, dashed, line width=1.0pt] (axis cs:45109, -280) -- (axis cs:45109, -135);
        \addplot[only marks, mark=*, mark size=2.6pt, purple] coordinates {(45109, -135)};
        \node[purple, fill=white, fill opacity=0.9, text opacity=1, inner sep=2pt, font=\small, anchor=north] 
            at (axis cs:45109, -143) {$f_{45} = 45,11\,\text{kHz}$};

        % Cruce f_180
        \draw[magenta!80!black, dashed, line width=1.0pt] (axis cs:84486, -280) -- (axis cs:84486, -180);
        \addplot[only marks, mark=square*, mark size=2.4pt, magenta!80!black] coordinates {(84486, -180)};
        \node[magenta!80!black, fill=white, fill opacity=0.9, text opacity=1, inner sep=2pt, font=\small, anchor=north] 
            at (axis cs:84486, -188) {$f_{180} = 84,49\,\text{kHz}$};
    \end{semilogxaxis}
\end{tikzpicture}
}
    \end{column}

    \begin{column}{0.40\textwidth}
      \begin{tarjeta}
        {\fontsize{8.5pt}{10pt}\selectfont \textbf{\color{NavyDark} Componentes y Desempeño:}}\\[2pt]
        {\fontsize{7pt}{8.5pt}\selectfont
        \begin{itemize}\setlength{\itemsep}{2pt}
          \item \textbf{Componentes ($C = \qty{100}{\pico\farad}$ C0G):}\\
          $R_2 = \mathbf{\qty{1.45}{\kilo\ohm}}$, $R_1 = \mathbf{\qty{5.76}{\kilo\ohm}}$ ($\qty{1}{\percent}$).
          \item \textbf{Ancho de banda preservado:}\\
          $BW \approx \mathbf{\qty{11.56}{\kilo\hertz}}$ ($\mathbf{99.7\%}$).
          \item \textbf{Frecuencia para $MF = \ang{45}$:}\\
          $f_{45} \approx \mathbf{\qty{45.11}{\kilo\hertz}}$ ($4.5\times$ más veloz).
          \item \textbf{Margen de Ganancia disponible:}\\
          $MG_{45} = \mathbf{8.96\ \unit{\decibel}} > \qty{6}{\decibel}$.
        \end{itemize}
        }
      \end{tarjeta}
    \end{column}
  \end{columns}
\end{frame}
